\documentclass[12pt, a4paper]{article}

% -- Encoding & fonts (XeLaTeX) ------------------------------
\usepackage{fontspec}
\setmainfont{Times New Roman}
\usepackage{unicode-math}
\setmathfont{XITS Math}

% -- Page geometry -------------------------------------------
\usepackage[
  a4paper,
  top    = 2.54cm,
  bottom = 2.54cm,
  left   = 2.54cm,
  right  = 2.54cm
]{geometry}

% -- Line spacing --------------------------------------------
\usepackage{setspace}
\doublespacing

% -- Running header / footer ---------------------------------
\usepackage{fancyhdr}
\setlength{\headheight}{14.5pt}
\pagestyle{fancy}
\fancyhf{}
\fancyhead[L]{\small\textit{Literature Review --- Mammographic Breast Cancer Detection}}
\fancyhead[R]{\small\thepage}
\renewcommand{\headrulewidth}{0.4pt}

% -- Section heading style -----------------------------------
\usepackage{titlesec}
\titleformat{\section}{\normalfont\bfseries\large}{\thesection}{1em}{}
\titleformat{\subsection}{\normalfont\bfseries}{\thesubsection}{1em}{}
\titlespacing*{\section}{0pt}{24pt}{6pt}
\titlespacing*{\subsection}{0pt}{18pt}{4pt}

% -- Bibliography ----------
\usepackage[
  backend   = biber,
  style     = ieee,
  sorting   = none
]{biblatex}
\addbibresource{litreview_refs_v4.bib}

% -- Hyperlinks ----------------------------------------------
\usepackage[
  colorlinks = true,
  linkcolor  = black,
  citecolor  = black,
  urlcolor   = blue
]{hyperref}
\usepackage{url}

% -- Misc ----------------------------------------------------
\usepackage{microtype}
\setlength{\parindent}{0.5in}
\setlength{\parskip}{0pt}

% -- Abstract environment ------------------------------------
\renewenvironment{abstract}{%
  \begin{center}\textbf{Abstract}\end{center}%
  \begin{quotation}%
}{\end{quotation}}

% ============================================================
\begin{document}

% -- Title page ----------------------------------------------
\begin{titlepage}
  \centering
  \vspace*{3cm}

  {\LARGE\bfseries Literature Review\par}
  \vspace{1.2cm}

  {\large\bfseries
    Mammographic Breast Cancer Detection with\\
    Conformal Uncertainty Quantification\par}
  \vspace{1.8cm}

  {\normalsize\itshape ACM 40960 \textperiodcentered\ Disease Modelling\par}
  \vspace{0.6cm}

  {\normalsize University College Dublin\\
  School of Mathematics and Statistics\par}
  \vspace{0.6cm}

  {\normalsize Supervised by Dr.\ Sarp Ak{\c{c}}ay\par}
  \vspace{0.6cm}

  {\normalsize Written by Darshan\ S and Rakshith\ K\ B\par}
  \vspace{0.6cm}
  {\normalsize June 2026\par}
\end{titlepage}

% -- Abstract ------------------------------------------------
\begin{abstract}
This review surveys the research landscape underpinning deep learning
system for breast cancer detection from digital mammography, organised
around the two methodological commitments that define the proposed work:
the fusion of bilateral information through an explicit asymmetry signal,
and the quantification of predictive uncertainty through distribution-free
conformal prediction. Rather than cataloguing prior systems individually,
The review develops the conceptual landscape from which these commitments
emerge, the architectural and clinical foundations of mammographic deep
learning, the progression from single-view through dual-view to bilateral
modelling, and the shift from heuristic confidence estimates toward formal
coverage guarantees and locates, within that landscape, the gap that the 
proposed system is designed to fill.
\end{abstract}

\newpage
\tableofcontents
\newpage

% ============================================================
\section{Introduction}

Breast cancer remains the most commonly diagnosed malignancy among women
worldwide, and mammographic screening is the primary tool through which it
is detected early enough for intervention to alter outcomes. The clinical
value of screening, however, is bounded by the limits of human reading:
sensitivity varies between radiologists, declines with fatigue across
high-volume screening sessions, and falls sharply when lesions are obscured
by radiographically dense tissue. It is against these limits that deep
learning has been positioned, with the most prominent systems now reporting
diagnostic performance comparable to, and in controlled studies exceeding,
that of experienced radiologists~\cite{McKinney2020, Wu2019}.

Progress at the level of headline performance has nonetheless left two
capabilities underdeveloped, and it is these two that the present work
takes as its subject. The first concerns how information from a screening
examination is combined. A standard mammographic study images each breast
from two angles, yet the diagnostic logic that radiologists apply to those
images --- the habitual comparison of left against right, in which a region
that departs from its contralateral counterpart draws immediate suspicion
--- is only partially reflected in systems that process views or breasts in
isolation. The second concerns how confidence is expressed. A screening
classifier that is wrong while appearing certain is more dangerous than one
that signals its own doubt, yet the uncertainty estimates accompanying most
deployed systems rest on heuristic foundations that carry no formal
guarantee.

This review establishes the conceptual context for both concerns. It traces
the architectural and clinical foundations on which mammographic deep
learning rests, follows the progression from single-view through dual-view
toward genuinely bilateral modelling, and examines the move from heuristic
confidence toward the distribution-free guarantees offered by conformal
prediction. The intent throughout is not to enumerate prior systems but to
develop the landscape they collectively define, and so to make visible the
specific gap --- the union of bilateral asymmetry fusion with conformal
uncertainty quantification --- that the proposed system occupies. Citations
follow the numeric convention; references are listed in full at the close.

% ============================================================
\section{Foundations of Mammographic Deep Learning}

The deep learning systems that now dominate mammographic analysis inherit
their representational machinery from a small lineage of convolutional
architectures, and the particular member of that lineage adopted here is
the densely connected network, in which each layer draws its input from the
concatenated outputs of all layers preceding it~\cite{Huang2017}. The
appeal of this connectivity pattern in a medical imaging setting is
practical rather than incidental: it sustains gradient flow through deep
stacks, encourages the reuse of features across layers, and reaches
competitive accuracy at a parameter budget modest enough to be fine-tuned
on datasets far smaller than those for which such networks were first
devised. That last property is decisive in mammography, where annotated
data is scarce relative to the natural-image corpora on which backbones are
pre-trained, and where the regularising effect of a compact, reuse-heavy
architecture translates directly into more reliable transfer.

The credibility of applying such architectures to screening at all was
established less by any single methodological advance than by demonstrations
of performance at clinical scale. The most visible of these evaluated a
system across screening populations drawn from more than one healthcare
system and reported reductions in both false positives and false negatives
relative to the standard of care, alongside a reading study in which the
system outperformed every radiologist enrolled~\cite{McKinney2020}. Read
narrowly, such results are a performance claim; read more broadly, they
fixed the expectations against which subsequent mammographic systems are
now measured, and they are the reference point any new architecture is
implicitly asked to justify itself against. A parallel line of academic
work, conducted at comparable scale but with its methodology and code
placed in the open, supplied what the proprietary demonstrations could
not~\cite{Wu2019}: a reproducible account of how a mammographic classifier
processes the multiple views of a screening study and combines them into a
case-level judgement. It is this account, rather than the headline
performance figures, that supplies the methodological baseline for the
present work.

Underlying both lines of work is a dataset dependency that shapes what is
practically buildable. The curated screening subset most widely used for
open mammographic research provides expert-annotated cases with pathology
labels, BI-RADS assessments, and lesion-level
delineations~\cite{Lee2017}, and its structure is consequential in ways
that are easy to underestimate. Cases are organised such that the two
breasts of a patient, and the two standard views of each breast, must be
reassembled by joining records on patient identity and laterality before
any bilateral relationship can be modelled at all; and because the
canonical train--test partition is defined at the level of the case rather
than the image, careless handling readily leaks one side of a patient into
training while the other is held out for test. The bilateral construction
on which the proposed system depends is therefore conditioned, from the
outset, on the organisation of this dataset, which makes its structure a
matter of architectural rather than merely administrative concern.

% ============================================================
\section{From Single Views to Bilateral Asymmetry}

The trajectory of mammographic modelling can be read as a steady widening of
the field of view that the classifier is permitted to consider at once. The
reproducible academic baseline already noted treats the four images of a
screening study together, weighting their contributions through learned
attention rather than collapsing them prematurely~\cite{Wu2019}, and the
immediate successor to that design made the relationship between the two
views of a single breast explicit, processing the craniocaudal and
mediolateral oblique projections jointly so that evidence visible in one
projection could inform interpretation of the other~\cite{Wu2020}. That
dual-view design is the closest architectural antecedent to the present
work, and the distinction between them is precisely the axis along which
the proposed system advances: where the antecedent fuses the two views of
the same breast, the proposed system additionally fuses across breasts,
treating the contralateral side not as a separate case to be classified but
as the reference against which asymmetry is measured.

The diagnostic motivation for that additional axis is well established in
the modelling literature, where the difference between corresponding left
and right representations has been used directly as a discriminative
signal. The most direct precedent computes an explicit asymmetry feature as
the absolute difference between the encoded representations of the two
breasts produced by a shared-weight encoder, and shows that supplying this
signal to the classifier improves discrimination over an otherwise
identical unilateral baseline~\cite{Nguyen2019}. The proposed system takes
this construction as its point of departure, retaining both the
shared-weight encoder and the difference signal \( \delta = |f_L - f_R| \),
and extends it with a gating mechanism that allows the magnitude of the
asymmetry to modulate how strongly each view representation contributes to
the final judgement.

That asymmetry can be exploited in more than one way is itself instructive.
A contrasting approach uses the bilateral relationship generatively,
synthesising a plausible healthy counterpart for a suspicious breast and
classifying by comparison against that counterfactual~\cite{Yang2021}. The
shared premise --- that the contralateral breast carries diagnostic
information no single image contains --- is exactly the premise the proposed
system adopts; the divergence lies in mechanism, and the divergence matters.
A generative counterpart introduces an auxiliary synthesis objective and the
attendant training complexity, and renders the basis of a decision harder to
inspect. By keeping asymmetry as a direct, differenced feature fused through
attention, the proposed system retains a representation whose contribution
can be visualised by standard gradient-based attribution, a property of
considerable value where a screening decision may ultimately require
justification to a clinician.

% ============================================================
\section{From Confidence to Coverage}

Quantifying how far a prediction may be trusted is, in a screening context,
not an adjunct to the classification problem but part of it. The estimates
of uncertainty that accompany most deployed systems are heuristic in origin
--- derived from the softmax response, from stochastic forward passes, or
from ensemble disagreement --- and while such estimates are often
empirically useful, they share a defining limitation: none carries a
guarantee about the rate at which the truth will actually fall within the
range the model declares. Conformal prediction occupies a different footing.
In its split formulation, a held-out calibration set is used to convert a
model's raw scores into prediction sets that satisfy a finite-sample
coverage guarantee, ensuring that for any chosen error level the true label
is contained with at least the corresponding probability, and doing so
without assumptions on the underlying data distribution~\cite{Angelopoulos2021}.
For a screening system, the consequence is concrete: rather than emitting a
single label and an uncalibrated confidence, the classifier can return a set
whose validity is guaranteed in a precise and verifiable sense, contracting
to a single class where the evidence is clear and expanding to admit both
possibilities where it is not.

The naive form of this construction, however, can yield sets larger than
necessary, diluting the practical value of the guarantee. A refinement
addresses this directly by adding a regularisation term that penalises the
inclusion of implausible classes, producing sets that remain valid in the
coverage sense while being substantially more economical in
size~\cite{Romano2020}. It is this refined construction that the proposed
system adopts, calibrating it on a dedicated split so that the prediction
sets it returns are both formally guaranteed and clinically legible: a
singleton when the model can commit, and a two-element set --- whose natural
reading is a recommendation to recall for further assessment --- when it
cannot.

% ============================================================
\section{Synthesis}

Surveyed together, these strands describe a landscape with a definite shape
and a definite vacancy. The architectural foundations supply a backbone
suited to data-scarce transfer~\cite{Huang2017}; the clinical and
reproducible baselines establish both the performance expectations of the
field and a concrete account of multi-view processing~\cite{McKinney2020,
Wu2019}; the dataset on which open research depends conditions what
bilateral modelling requires in practice~\cite{Lee2017}. Along the modelling
axis, the progression runs from multi-view~\cite{Wu2019}, through
dual-view~\cite{Wu2020}, to the explicit use of bilateral asymmetry as a
differenced feature~\cite{Nguyen2019} or as a generative
prior~\cite{Yang2021}. Along the uncertainty axis, the movement is from
heuristic confidence toward distribution-free guarantees~\cite{Angelopoulos2021},
made practical by size-regularised prediction sets~\cite{Romano2020}.

What no single line within this landscape provides is their union. The
bilateral systems quantify no uncertainty in any guaranteed sense; the
conformal methods are developed without reference to the bilateral structure
of mammography; and the closest architectural antecedent stops at the
boundary between the two views of one breast, short of the contralateral
comparison that defines clinical reading. The proposed system is situated
precisely at that intersection, joining an explicit, attention-gated
bilateral asymmetry signal to a distribution-free conformal guarantee within
a single architecture trained on openly available data --- a combination
that the surveyed literature motivates at every point but nowhere assembles.

% ============================================================
\section{Conclusion}

The foregoing has developed, rather than enumerated, the context for the
proposed system. From the architectural lineage it draws a backbone matched
to the data regime of mammography; from the clinical and academic baselines
it inherits both a standard of performance and a reproducible model of
multi-view processing; from the asymmetry literature it takes the differenced
bilateral signal it extends, and the generative alternative it deliberately
declines; and from the conformal literature it adopts the guarantee, and the
refinement of that guarantee, on which its treatment of uncertainty rests.
The scope of the review has been kept deliberately narrow, admitting only
those works that bear directly on the two commitments the system is built
around, so that the gap it occupies --- the conjunction of bilateral
asymmetry fusion with conformal uncertainty quantification --- stands clear
of the surrounding field.

% ============================================================
\newpage
\printbibliography[title={References}]

\end{document}
